% !PRINTED none
% Abschnitt-opening page: no running head/page number printed on this recto.
\begin{center}
\gesperrt{Zwölfter Abschnitt.}

\textbf{Die Doppeltangenten einer Curve vierter Ordnung.}

\rule{0.3\linewidth}{0.4pt}

§. 95.

\gesperrt{Anzahl der Doppeltangenten einer Curve\\
vierter Ordnung.}
\end{center}

Ein geometrisches Problem, das wegen seiner mannigfachen Beziehungen zu anderen Gebieten von besonderer Bedeutung ist, was zugleich in ähnlicher Weise, wie das Problem der Wendepunkte der Curven dritter Ordnung, merkwürdige algebraische Verhältnisse bietet, ist das der Doppeltangenten der Curven vierter Ordnung.

Unter einer \gesperrt{Doppeltangente} einer Curve versteht man, wie schon im \S.\,86 bemerkt ist, eine gerade Linie, die die Curve in zwei verschiedenen Punkten berührt. Bei Curven von niedrigerer als der vierten Ordnung können Doppeltangenten nicht auftreten.

Bei den Curven vierter Ordnung hat eine Doppeltangente ausser den Berührungspunkten keinen Punkt mit der Curve gemein. In besonderen Fällen können die beiden Berührungspunkte auch zusammenfallen. Dann haben wir Linien mit vierpunktiger Berührung.

Die Bestimmung der Doppeltangenten wird als algebraisches Problem von einer gewissen algebraischen Gleichung abhängen, deren Grad gleich der Anzahl der Doppeltangenten ist, und die erste Frage ist die nach dem Grade dieser Gleichung, also nach der Anzahl der Doppeltangenten. Wir beschränken uns hier auf die Betrachtung von Curven vierter Ordnung, obwohl der Weg, den wir gehen, auch auf Curven höherer Ordnung anwendbar
%CONT>>
